\documentclass{article}


\usepackage{geometry}
\geometry{
    textwidth=15cm,  % 增加文本区域的宽度
    margin=2.5cm    % 调整页边距
}


\usepackage{microtype}
\usepackage{amsmath}
\usepackage{amssymb}
\usepackage{graphicx}
\usepackage{listings}
\usepackage{epstopdf}
\graphicspath{ {images/} }


\title{PHY2650 Homework1}
\author{Chen Yilun(123090048)}


\begin{document}

\maketitle


All codes please refer to the attached python files!

All codes please refer to the attached python files!

All codes please refer to the attached python files!



\section{Q1}

\begin{figure}[htbp]
    \centering
    \includegraphics[width=10cm]{images/hw1q1.jpg}
\end{figure}

\section{Q2}
\subsection{Code}
\begin{lstlisting}[language=Python]
import numpy as np  
import matplotlib.pyplot as plt  # Import the pyplot module
import math  


def f(x):
    """
    Define the function f(x) 
    :param x: The input independent variable
    :return: The result of the function calculation
    """
    return x**6 + 0.1 * np.log(np.abs(1 + 3 * (1 - x)))


numpoints = 100  # Define the number of points used to plot the function


x = np.linspace(0.5, 1.5, numpoints)  # Generate numpoints in the interval 


y = f(x)  


plt.xlabel("x")  # Set the label of the x-axis
plt.ylabel("f(x)")  # Set the label of the y-axis


plt.plot(x, y)  # Plot the function graph
plt.show()  # Display the function graph


k = 1.3  # Initialize the variable k to 1.3
while k <= 1.4:  
    print(k, f(k))  
    k += 0.01  # Increase the value of k by 0.01
\end{lstlisting}



\subsection{Plot(Figure 1 and Figure 2)}

\begin{figure}[htbp]
    \centering
    \includegraphics[width=7cm]{images/HW1Q2GRID.png}
    \caption{plot of f(x)}
    \label{plot of f(x)}
\end{figure}

\begin{figure}[h]
    \centering
    \includegraphics[width=7cm]{images/HW1Q2GRIDFIN.png}
    \caption{plot of f(x) with finer grid}
    \label{plot of f(x) with finer grid}
\end{figure}


The ``dip'' occurs at \(x = \frac{4}{3}\).

The ``dip'' of the function is mainly due to the logarithmic part. 

When \(x\) approaches \(\frac{4}{3}\), 

the logarithmic function \(y = 0.1\log(|4 - 3x|)\) 

decreases rapidly because the logarithmic function approaches negative infinity 

as it approaches zero from the positive direction. 

At this time, 

the growth of the polynomial part (\(x^{6}\)) 

is not fast enough to offset the decreasing trend of the logarithmic part.

When \(x\) exceeds \(\frac{4}{3}\), 

the growth of the polynomial part dominates, 

and the function starts to rise again.



\section{Q3}
\begin{lstlisting}
def searchSquareNum(n):
    numList = []  # Initialize an empty list to store square numbers
    
    try:
        n = int(n)  
    except:
        print('Please input only an integer')  
        #send error message if the input n is not an integer
    
    k = 1
    m = 1
    
    while True:
        k = m ** 2  # Calculate the square of m
        if k >= n:  # Check if square is greater than or equal to n
            break
        numList.append(k)  # Append the square to numList
        m += 1  # Increment m
    
    print(f'The square numbers up to {n} are {numList}')  
    # Print the result in format
    return numList


searchSquareNum(1000)

\end{lstlisting}

\section{Q4}

\begin{lstlisting}
import math


def taylorSineTerm(x, N):
    # Calculate the Nth term of the Taylor series for sine
    return ((-1) ** N) * (x ** (2 * N + 1)) / math.factorial(2 * N + 1)


def taylorExpansionOfSin(x):
    sum = 0  # Initialize the sum of the Taylor series
    N = 0  # Initialize the term index
    if x == 0:
        print(0)
        print(f'the summation ends at N = {N}')
        return 0
    while abs(taylorSineTerm(x, N) * (10 ** 7)) >= abs(sum):
        # Add the Nth term to the sum if it's significant
        sum += taylorSineTerm(x, N)
        N += 1
        # print(taylorSineTerm(x, N))
    print(sum)
    print(f'the summation ends at N = {N}')
    return sum


def difference(x):
    # Calculate the difference between the Taylor expansion 
    # and the actual value
    diff = abs(taylorExpansionOfSin(x) - math.sin(x))
    print(f'the difference is {diff}')
    return diff


# Try x = 0.5
taylorExpansionOfSin(0.5)
print()
difference(0.5)
\end{lstlisting}




\section{Q5}
Here are the values and error when x is from 1 to 100:
\begin{lstlisting}
x = 1, N = 5, sum = 0.8414710097001764, error = 2.4892279903987458e-08
x = 2, N = 7, sum = 0.9092974515196738, error = 2.469399207338796e-08
x = 3, N = 9, sum = 0.1411200174132425, error = 9.353375296505817e-09
x = 4, N = 11, sum = -0.7568024926569087, error = 2.6510195150208915e-09
x = 5, N = 12, sum = -0.9589242932128198, error = 1.854968134651358e-08
x = 6, N = 14, sum = -0.27941550221018424, error = 4.011258381009242e-09
x = 7, N = 15, sum = 0.6569866170512744, error = 1.8332485307581692e-08
x = 8, N = 16, sum = 0.9893581773748686, error = 6.924851314593639e-08
x = 9, N = 18, sum = 0.4121184712778701, error = 1.3963886491552557e-08
x = 10, N = 19, sum = -0.5440210646963509, error = 4.619301885444571e-08
x = 10, N = 19, sum = -0.5440210646963509, error = 4.619301885444571e-08
x = 11, N = 21, sum = -0.9999901971573234, error = 9.393380073774438e-09
x = 12, N = 22, sum = -0.536572946655428, error = 2.8654993045051924e-08
x = 13, N = 24, sum = 0.4201670309237957, error = 5.902845212624186e-09
x = 14, N = 25, sum = 0.990607372754533, error = 1.7059662682505916e-08
x = 15, N = 26, sum = 0.6502877933350776, error = 4.6822039245775215e-08
x = 16, N = 28, sum = -0.28790332666974433, error = 1.0004679029673724e-08
x = 17, N = 29, sum = -0.9613974655323653, error = 2.6347191539599635e-08
x = 18, N = 30, sum = -0.7509873140223561, error = 6.725067991375511e-08
x = 19, N = 33, sum = 0.14987721106093257, error = 1.3979802315766676e-09
x = 20, N = 33, sum = 0.9129452856899325, error = 3.496230482902263e-08
x = 21, N = 35, sum = 0.8366556502462718, error = 1.1710215686022707e-08
x = 22, N = 38, sum = -0.00885130517331938, error = 4.117084496241263e-09
x = 23, N = 37, sum = -0.8462204239602107, error = 1.9785040050734892e-08
x = 24, N = 39, sum = -0.9055783071852938, error = 5.482133003109624e-08
x = 25, N = 41, sum = -0.1323523332756863, error = 5.831779132658976e-07
x = 26, N = 41, sum = 0.7625585677098587, error = 1.1723025605814286e-07
x = 27, N = 43, sum = 0.9563779283175255, error = 1.999913022521227e-06
x = 28, N = 45, sum = 0.2709060649291043, error = 2.766212352511843e-07
x = 29, N = 46, sum = -0.6636555417814031, error = 2.165756843552913e-05
x = 30, N = 47, sum = -0.9882783116573879, error = 0.0002466875645260602
x = 31, N = 49, sum = -0.4039467007790409, error = 9.094454402414076e-05
x = 32, N = 50, sum = 0.5520087616982938, error = 0.00058208045660324
x = 33, N = 51, sum = 0.9980331969001163, error = 0.0018786632071509102
x = 34, N = 53, sum = 0.5193320744349768, error = 0.009750611685047095
x = 35, N = 54, sum = -0.43441292980413043, error = 0.006230260307979407
x = 36, N = 55, sum = -0.9650706547132785, error = 0.026708198729837318
x = 37, N = 56, sum = -0.7565588938026231, error = 0.11302076044562359
x = 38, N = 58, sum = 0.3960663661733253, error = 0.09969778746393998
x = 39, N = 60, sum = 0.3236822919430559, error = 0.6401130943410318
x = 40, N = 60, sum = -3.2822092133102636, error = 4.027322373789612
x = 41, N = 61, sum = -2.812198282315549, error = 2.6535756135108404
x = 42, N = 62, sum = -11.189445057817396, error = 10.272923509901762
x = 43, N = 63, sum = -67.27444478273009, error = 66.44267004010149
x = 44, N = 64, sum = -69.08329863432868, error = 69.1010005594341
x = 45, N = 66, sum = 37.93452239839774, error = 37.08361887386362
x = 46, N = 66, sum = 256.61056770958965, error = 255.70877936194083
x = 47, N = 67, sum = -762.0169304979033, error = 762.1405036206486
x = 48, N = 67, sum = 3458.5643596590594, error = 3459.332614320383
x = 49, N = 68, sum = -19811.728971354867, error = 19810.775218702107
x = 50, N = 69, sum = -25906.59440991337, error = 25906.332035059666
x = 51, N = 70, sum = -172896.6928918742, error = 172897.36312105006
x = 52, N = 71, sum = 441930.06314069906, error = 441929.076513107
x = 53, N = 72, sum = 798689.1039335495, error = 798688.7080083993
x = 54, N = 74, sum = 201091.35277015856, error = 201091.9115592074
x = 55, N = 74, sum = 774825.1672000011, error = 774826.1669551744
x = 56, N = 74, sum = 19562797.674335428, error = 19562798.19588643
x = 57, N = 76, sum = -13733634.27434589, error = 13733634.710510645
x = 58, N = 76, sum = -46164506.101266704, error = 46164507.09413935
x = 59, N = 77, sum = -394867750.18983245, error = 394867750.82657045
x = 60, N = 77, sum = -1040182759.7300932, error = 1040182759.4252826
x = 61, N = 78, sum = -1902532877.5632563, error = 1902532876.5971384
x = 62, N = 79, sum = 13427775720.745657, error = 13427775721.484838
x = 63, N = 80, sum = 8053116623.358149, error = 8053116623.190793
x = 64, N = 81, sum = -63969027487.9773, error = 63969027488.89733
x = 65, N = 81, sum = -194932344570.60422, error = 194932344571.43106
x = 66, N = 83, sum = -249359088183.4633, error = 249359088183.43674
x = 67, N = 83, sum = 1470799996627.6624, error = 1470799996628.5178
x = 68, N = 84, sum = 1736627568267.4587, error = 1736627568268.3567
x = 69, N = 85, sum = -7179210615847.565, error = 7179210615847.45
x = 70, N = 87, sum = -3164386588246.2295, error = 3164386588247.0034
x = 71, N = 89, sum = -519990991575.5949, error = 519990991576.54596
x = 72, N = 88, sum = 18817914726329.14, error = 18817914726328.887
x = 73, N = 90, sum = 21223128584701.527, error = 21223128584702.203
x = 74, N = 89, sum = 389913248090217.0, error = 389913248090218.0
x = 75, N = 90, sum = 1162792272048649.2, error = 1162792272048649.8
x = 76, N = 90, sum = -8205109641752168.0, error = 8205109641752169.0
x = 77, N = 91, sum = 2.488303381721387e+16, error = 2.488303381721387e+16
x = 78, N = 92, sum = -2.905294170961408e+16, error = 2.905294170961408e+16
x = 79, N = 93, sum = 8.430036165125453e+16, error = 8.430036165125453e+16
x = 80, N = 94, sum = 9.487840541325083e+16, error = 9.487840541325083e+16
x = 81, N = 95, sum = 8.817233029809887e+17, error = 8.817233029809887e+17
x = 82, N = 96, sum = -1.1187322291849094e+18, error = 1.1187322291849094e+18
x = 83, N = 96, sum = 4.968939622441953e+18, error = 4.968939622441953e+18
x = 84, N = 97, sum = -2.0025445713603867e+19, error = 2.0025445713603867e+19
x = 85, N = 98, sum = -5.408322702569982e+19, error = 5.408322702569982e+19
x = 86, N = 98, sum = -1.4483709569991347e+20, error = 1.4483709569991347e+20
x = 87, N = 99, sum = 2.719619975211302e+20, error = 2.719619975211302e+20
x = 88, N = 100, sum = 6.83289687346446e+20, error = 6.83289687346446e+20
x = 89, N = 100, sum = -4.804950979144479e+21, error = 4.804950979144479e+21
x = 90, N = 102, sum = -7.005429166281726e+21, error = 7.005429166281726e+21
x = 91, N = 105, sum = 1.496404659458053e+20, error = 1.496404659458053e+20
x = 92, N = 103, sum = -5.850410323941645e+22, error = 5.850410323941645e+22
x = 93, N = 104, sum = 9.658404941266049e+22, error = 9.658404941266049e+22
x = 94, N = 106, sum = 4.138936437904746e+22, error = 4.138936437904746e+22
x = 95, N = 107, sum = -7.903539633668151e+22, error = 7.903539633668151e+22
x = 96, N = 107, sum = -1.5972826674735094e+24, error = 1.5972826674735094e+24
x = 97, N = 107, sum = 3.3126774191974173e+24, error = 3.3126774191974173e+24
x = 98, N = 107, sum = -3.335885370084312e+25, error = 3.335885370084312e+25
x = 99, N = 108, sum = -7.349906874266141e+25, error = 7.349906874266141e+25
x = 100, N = 110, sum = 3.6809885386809145e+25, error = 3.6809885386809145e+25

\end{lstlisting}

From above we can observe that 

when x <= 24, the relative error is of $10^{-9}$ or $10^{-8}$;

when x>=28, the relative error started to get larger than $10^{-7}$; 

when x>=32, the relative error is above $10{-3}$, where the accuracy become worse;

when x>=39, the relative error is comparable or larger than 1 and the accuracy is totally lost.

When x>=40, the series no longer converges since the output of the sum is 

larger than or much larger than 1.

If x continues to grow, the sum and error are ridiculously large.







\section{Q6}
\subsection{Code}
\begin{lstlisting}
import math
import numpy as np
import matplotlib.pyplot as plt

# Function to calculate the n-th term of the Taylor series for sin(x)
def taylorSineTerm(x, n):
    return ((-1) ** n) * (x ** (2 * n + 1)) / math.factorial(2 * n + 1)

# Function to calculate the error of the Taylor 
series approximation of sin(x) up to N terms
def taylorExpansionOfSin(x,N):
    sum_result = 0
    n = 0
    while n <= N:
        term = taylorSineTerm(x, n)
        sum_result += term
        n += 1
    error = abs(sum_result - math.sin(x))
    return error


# List of x values to calculate errors for
x_values = [0.5, 1, 2, 5, 10]

for x in x_values:
    N_values = []
    Error_values = []
    max_N = 20
    
    # Calculate errors for different numbers of terms (N)
    for N in range(1,max_N):
        error = taylorExpansionOfSin(x,N)
        N_values.append(N)
        Error_values.append(error)
    
    N_values = np.array(N_values)
    Error_values = np.array(Error_values)
    
    plt.plot(N_values, Error_values, label=f'x = {x}')

plt.xlabel('N')
plt.ylabel('Error')

plt.xticks(range(0, max_N + 1, 1))   # Set x-axis tick marks
plt.ylim(0, 0.3)

plt.title('Error vs N for different x values')  # Set y-axis limits
plt.legend()  # Show legend for x with different colors

plt.show()




\end{lstlisting}

\subsection{Plot(See Figure 3 next page)}
\begin{figure}[htbp]
    \centering
    \includegraphics[width=7cm]{images/hw1q6.png}
    \caption{plot of error vs. N}
    \label{Error vs. N for Taylor Sine Expansion}
\end{figure}


\section{Q7}
\subsection{Code}
\begin{lstlisting}
import math
import numpy as np
import matplotlib.pyplot as plt

# Initialize variables
D = {} #The dictionary stores N as keys and results as values
Errors = {}  #The dictionary stores N as keys and errors as values

# Define the function xloop to calculate x for different N values
def xloop():
    # Loop through N from 10 to 100
    for N in range(10, 101):
        # Reset x to 1/2 at the start of each N iteration
        x = 1/2
        # Perform N nested calculations
        for _ in range(1, N + 1):
            x = 1 / (2 + x)
        # print(f'When N = {N}, x = {x}')
        # Store the current N and x in dictionary D
        D[N] = 1 + x
        Errors[N] = abs((1+x) - math.sqrt(2))


xloop()



# Extract N values and corresponding results
N_values = list(D.keys())
result_values = list(D.values())
# Transform into arrays
N_values = np.array(N_values)
result_values = np.array(result_values)



# Plot the results as a function of N
plt.figure(figsize=(10, 6))
plt.plot(N_values, result_values)
plt.title('Results vs N')
plt.xlabel('N')
plt.ylabel('Result (x)')
plt.grid(True)


# Show the plot
plt.show()




# Extract N values and corresponding errors
N_values = list(Errors.keys())
error_values = list(Errors.values())
# Transform into arrays
N_values = np.array(N_values)
error_values = np.array(error_values)



# Plot the errors as a function of N
plt.figure(figsize=(10, 6))
plt.plot(N_values, error_values)
plt.title('Errors vs N')
plt.xlabel('N')
plt.ylabel('Error (x)')
plt.grid(True)


# Show the plot
plt.show()
\end{lstlisting}



\subsection{Plot1(Figure 4)}
\begin{figure}[htbp]
    \centering
    \includegraphics[width=7cm]{images/hw1q7（rvn）.png}
    \caption{plot of result vs. N}
    \label{Result vs. N for Taylor Sine Expansion}
\end{figure}



\subsection{Plot2(See Figure 5 next page)}
\begin{figure}[htbp]
    \centering
    \includegraphics[width=7cm]{images/hw1q7（evn）.png}
    \caption{plot of error vs. N}
    \label{Error vs. N for Taylor Sine Expansion}
\end{figure}




\section{Q8}
\subsection{Code}
\begin{lstlisting}
import math

def singleContribution(i, j, k):
    '''Define a function to calculate the contribution of a single atom 
    to the Madelung constant'''
    if i == 0 and j == 0 and k == 0:
        return 0
    
    absContribution = math.sqrt(i**2 + j**2 + k**2) 
    #calculate the absolute value of the contribution
    
    if (i + j + k) % 2 == 0: #if the sum of i, j, k 
    is even, 
    the contribution is positive
        return 1 / absContribution
    else:                    #if the sum of i, j, k is odd, 
    the contribution is negative
        return -1 / absContribution
    

def getMadelung(L):
    ''''Define a function to calculate the Madelung constant
    for a range of L'''
    sum = 0
    for i in range(-L, L+1):
        for j in range(-L, L+1):
            for k in range(-L, L+1):
                sum += singleContribution(i, j, k)  
                #sum the contribution of each atom
    return sum

L = 200 #set the range of L
print(getMadelung(L))
#the result is -1.7446850421707383 when L = 200
#using C++ can get the result faster than python
\end{lstlisting}

\subsection{Result}
The result is -1.7446850421707383 or the absolute value 1.7446850421707383

\end{document}

